\section*{Блок 4. Ожидания функций от независимых величин и смешанные распределения}
\addcontentsline{toc}{section}{Блок 4. Ожидания функций от независимых величин и смешанные распределения}

\subsection*{1. Семантический анализ и принцип итерированного ожидания}
\addcontentsline{toc}{subsection}{1. Семантический анализ и принцип итерированного ожидания}

Специфика данного вычислительного блока заключается в отсутствии прямых аналогов в семинарских занятиях. Задача №4 из контрольной работы проверяет понимание конструкции интеграла Лебега-Стилтьеса и умение работать со смешанными (кусочно-заданными с разрывами) распределениями на практическом уровне. 

Основная проблема формулируется следующим образом: даны две независимые случайные величины $X$ и $Y$. Требуется вычислить математическое ожидание нелинейной функции $\E[g(X, Y)]$. При этом распределение $X$ имеет как участки непрерывного роста, так и точки разрыва первого рода (скачки), а $Y$ абсолютно непрерывна.

Прямое применение двумерного интеграла Лебега-Стилтьеса в таких задачах громоздко. Оптимальный аналитический подход базируется на \textbf{свойстве башни (телескопическом свойстве)} условного математического ожидания и независимости величин. Поскольку $X \indep Y$, условное математическое ожидание $\E[g(X, Y) \mid X]$ сводится к вычислению функции от одной переменной. Введем детерминированную функцию $\psi(x)$:
\[
    \psi(x) \coloneqq \E[g(x, Y)] = \int_{\mathbb{R}} g(x, y) p_Y(y) \, dy.
\]
Тогда искомая величина вычисляется как безусловное математическое ожидание функции от $X$:
\[
    \E[g(X, Y)] = \E[\E[g(X, Y) \mid X]] = \E[\psi(X)] = \int_{\mathbb{R}} \psi(x) \, dF_X(x).
\]
Данная редукция позволяет свести двумерную задачу к двум последовательным одномерным интегралам, где специфика смешанного распределения учитывается только на последнем этапе.

\subsection*{2. Теоретический базис: Декомпозиция Лебега и аппарат обобщенных функций}
\addcontentsline{toc}{subsection}{2. Теоретический базис: Декомпозиция Лебега и аппарат обобщенных функций}

Согласно теореме Лебега, любая функция распределения $F_X(x)$ может быть единственным образом представлена в виде:
\[
    F_X(x) = F_{ac}(x) + F_d(x) + F_{sing}(x),
\]
где в рамках учебного курса сингулярная компонента $F_{sing}(x) \equiv 0$, $F_{ac}(x)$ — абсолютно непрерывная часть, а $F_d(x)$ — дискретная функция скачков.

Вместо классического разбиения интеграла на две независимые части (сумму и интеграл Римана), вычислительно эффективнее использовать аппарат обобщенных функций.

\begin{HackBox}[Универсальная плотность через дельта-функцию Дирака]
Пусть $F_X(x)$ дифференцируема почти всюду, за исключением счетного множества точек разрыва первого рода $D = \{x_1, x_2, \dots, x_k\}$. 
Определим обобщенную плотность $p^*_X(x)$:
\[
    p^*_X(x) = F_{ac}'(x) + \sum_{x_k \in D} \P(X = x_k) \cdot \delta(x - x_k),
\]
где $p_k = \P(X = x_k) = F_X(x_k) - F_X(x_k - 0)$ — величина скачка функции распределения.

Использование фильтрующего свойства дельта-функции $\int_{\mathbb{R}} \psi(x)\delta(x - x_k)dx = \psi(x_k)$ сводит итоговое вычисление к единому интегралу:
\[
    \E[\psi(X)] = \int_{\mathbb{R}} \psi(x) p^*_X(x) \, dx = \int_{\mathbb{R} \setminus D} \psi(x) F_{ac}'(x) \, dx + \sum_{x_k \in D} \psi(x_k) p_k.
\]
Этот метод минимизирует риск потери одной из компонент (дискретной или непрерывной) при оформлении решения.
\end{HackBox}

\subsection*{3. Универсальный алгоритм решения}
\addcontentsline{toc}{subsection}{3. Универсальный алгоритм решения}

\begin{MethodBox}[Алгоритм вычисления ожидания для смешанного распределения]
Рассматривается задача вычисления $\E[g(X, Y)]$, где $F_X(x)$ имеет скачки, а $Y$ непрерывна.

\begin{enumerate}[label=\textbf{\arabic*.}]
    \item \textbf{Интегрирование по непрерывной величине.} Зафиксировать $X = x$ как параметр. Вычислить внутреннее математическое ожидание $\psi(x) = \E[g(x, Y)]$ интегрированием по плотности $p_Y(y)$.
    \item \textbf{Локализация атомов (скачков).} Исследовать $F_X(x)$ на точки разрыва. Найти координаты разрывов $x_k$, в которых существует ненулевой предел слева, не равный значению функции. Вычислить вероятностные массы $p_k = F_X(x_k) - \lim_{\varepsilon \to 0^+} F_X(x_k - \varepsilon)$.
    \item \textbf{Дифференцирование непрерывной части.} На интервалах непрерывности найти производную $p_{ac}(x) = F_X'(x)$. В точках разрыва положить производную равной нулю.
    \item \textbf{Синтез.} Применить формулу полного математического ожидания (интегрирование по обобщенной плотности):
    \[
        \E[g(X, Y)] = \sum_{k} \psi(x_k) p_k + \int_{-\infty}^{+\infty} \psi(x) p_{ac}(x) \, dx.
    \]
\end{enumerate}
\end{MethodBox}

\vspace{1em}

\subsection*{4. Оптимизация порядка интегрирования (Теорема Тонелли)}
\addcontentsline{toc}{subsection}{4. Оптимизация порядка интегрирования (Теорема Тонелли)}

В задачах вида $\E[Y^X]$ порядок интегрирования критически важен для минимизации временных затрат. Теорема Тонелли для неотрицательных измеримых функций гарантирует независимость результата от порядка интегрирования, однако аналитическая сложность путей различается радикально.

\begin{HackBox}[Аналитический срез для показательных функций]
Если функция имеет вид $g(x, y) = y^x$ и $Y \sim U(a, b)$, интегрирование по $Y$ при фиксированном параметре $X=x$ сводится к интегрированию степенной функции:
\[
    \psi(x) = \E[Y^x] = \int_a^b y^x \frac{1}{b-a} \, dy = \frac{1}{b-a} \cdot \frac{b^{x+1} - a^{x+1}}{x+1} \quad (\text{для } x \neq -1).
\]
Дальнейшее интегрирование $\psi(x)$ по мере Стилтьеса $dF_X(x)$ требует лишь взятия интегралов от показательных функций $b^x$ и $a^x$. 
Альтернативный путь (попытка интегрировать $y^x = e^{x \ln y}$ по мере $X$ в первую очередь) приводит к необходимости работать с трансцендентными функциями вида $e^{C \ln y}$ по смешанной мере, что многократно усложняет преобразования и повышает вероятность алгебраической ошибки.
\end{HackBox}

\subsection*{5. Опасные места и эвристики самопроверки}
\addcontentsline{toc}{subsection}{5. Опасные места и эвристики самопроверки}

\begin{WarningBox}[Ложная нормировка непрерывной компоненты]
Фундаментальная ошибка при работе со смешанными распределениями — попытка воспринимать непрерывную часть $p_{ac}(x) = F_X'(x)$ как полноценную плотность вероятности. 
Интеграл от производной непрерывной части по всей числовой прямой \textbf{строго меньше единицы}:
\[
    \int_{-\infty}^{+\infty} p_{ac}(x) \, dx = 1 - \sum_{k} \P(X = x_k).
\]
Применение стандартных табличных формул математического ожидания к $p_{ac}(x)$ в отрыве от дискретной части недопустимо. Функция $\psi(x)p_{ac}(x)$ должна интегрироваться строго по определению в пределах носителя непрерывной компоненты.
\end{WarningBox}

\begin{HackBox}[Алгоритмы граничного и логического контроля]
Для мгновенной верификации полученного ответа $\E[g(X, Y)]$ применяются следующие аналитические эвристики:
\begin{enumerate}
    \item \textbf{Закон сохранения вероятностной массы.} Перед вычислением финального интеграла вычислите сумму скачков и интеграл от производной непрерывной части: $\sum p_k + \int p_{ac}(x) dx$. Сумма обязана быть строго равна $1$. Это базовый критерий корректности декомпозиции.
    \item \textbf{Тест константы (Вырождение).} Если подставить в итоговую формулу вырожденное распределение $Y \equiv 1$ (что достигается предельным переходом $Y \sim U(1-\varepsilon, 1+\varepsilon)$ при $\varepsilon \to 0$), то для $g(X,Y) = Y^X$ должно строго выполняться $\E[1^X] = 1$. Если аналитический ответ при $b \to 1, a \to 1$ не стремится к $1$, допущена ошибка интегрирования.
    \item \textbf{Оценка монотонности.} Если носитель $Y$ лежит в интервале $(1, \infty)$, а носитель $X$ строго положителен ($X \ge 0$ почти наверное), то $Y^X \ge 1$ с вероятностью 1. Следовательно, $\E[Y^X] \ge 1$. Получение значения меньше единицы сигнализирует о потере слагаемого или ошибке в знаке.
\end{enumerate}
\end{HackBox}
